\documentclass{article}
\usepackage[a4paper, margin=1in]{geometry}
\usepackage{amsmath, amssymb}
\usepackage{enumitem}
\usepackage{hyperref}

\title{}
\date{}

\begin{document}

\maketitle

\section*{Domination in Unit Disk Graphs}
Nieberg and Hurink~\cite{nieberg2006ptas} proposed a polynomial-time approximation scheme (PTAS) for the Minimum Dominating Set problem in Unit Disk Graphs (UDGs) that does not rely on the geometric representation of the graph, achieving a $(1+\varepsilon)$-approximation in time $n^{O(1/\varepsilon \log (1/\varepsilon))}$, and also providing a robustness feature by either returning a valid solution or a certificate that the input is not a UDG.


\section*{Semi-total Domination in Unit Disk Graphs}


\textbf{Authors:} Sasmita Rout, Gautam Kumar Das

\textbf{Published in:} CALDAM 2024 (LNCS 14508)

\textbf{DOI:} \href{https://doi.org/10.1007/978-3-031-52213-0_9}{10.1007/978-3-031-52213-0\_9}

\vspace{1em}

\section*{1. Problem Studied}
\begin{itemize}[leftmargin=1.5em]
    \item The paper studies the \textbf{Semi-total Dominating Set (STDS)} problem in Unit Disk Graphs (UDGs).
    \item A set \( D_{t2} \subseteq V \) is a STDS if:
    \begin{enumerate}
        \item \( D_{t2} \) is a dominating set.
        \item For every \( u \in D_{t2} \), there exists \( v \in D_{t2} \) such that \( \text{dist}(u,v) \leq 2 \).
    \end{enumerate}
\end{itemize}

\section*{2. Motivation and Applications}
\begin{itemize}[leftmargin=1.5em]
    \item STDS is a natural variant lying between classical dominating sets and total dominating sets.
    \item Relevant in wireless networks and sensor placements where both coverage and minimal connectivity are required.
\end{itemize}

\section*{3. Main Contributions}
\begin{itemize}[leftmargin=1.5em]
    \item Proved that the STDS problem is \textbf{NP-complete} in UDGs via a reduction from Vertex Cover on planar graphs.
    \item Proposed a \textbf{6-factor approximation algorithm} for STDS in UDGs.
    \item Introduced a \textbf{2 + ln(D+1)} approximation for general graphs, improving over the previous bound of \textbf{2 + 3ln(D+1)}.
    \item Showed that a PTAS for STDS in UDGs exists with an approximation factor of \( 2 + \varepsilon \).
\end{itemize}

\section*{4. Techniques Used}
\begin{itemize}[leftmargin=1.5em]
    \item \textbf{NP-Completeness:} Used grid embedding and geometric gadget constructions to simulate vertex cover behavior.
    \item \textbf{Approximation Algorithm:}
    \begin{enumerate}
        \item Compute a maximal independent set \( D \).
        \item For each vertex in \( D \) lacking a neighbor within distance 2 in \( D \), add a neighbor from \( N(u) \) to set \( T \).
        \item Return \( D \cup T \) as the semi-total dominating set.
    \end{enumerate}
\end{itemize}

\section*{5. Results Summary}

\begin{tabular}{|l|l|l|l|}
\hline
\textbf{Graph Type} & \textbf{Approach} & \textbf{Approx. Factor} & \textbf{Time Complexity} \\
\hline
UDG & Greedy + Local Fix & 6 & \( O(nk) \) \\
UDG (PTAS) & Local Search + Packing & \( 2 + \varepsilon \) & \( n^{O(1/\varepsilon \log \varepsilon)} \) \\
General Graph & Set Cover-Based & \( 2 + \ln(D + 1) \) & --- \\
\hline
\end{tabular}

\vspace{1em}

\section*{6. Comparison with Prior Work}
\begin{itemize}[leftmargin=1.5em]
    \item Improves over the \( 2 + 3 \ln(D+1) \) result from Henning and Pandey (2019).
    \item The first work to give a dedicated algorithm and complexity result for STDS in UDGs.
    \item Relates to: \( \gamma(G) \leq \gamma_{t2}(G) \leq \gamma_t(G) \).
\end{itemize}

\section*{7. Limitations / Assumptions}
\begin{itemize}[leftmargin=1.5em]
    \item Assumes no isolated vertices.
    \item Requires geometric embedding (disk centers must be known).
    \item Maximal independent set construction may not be efficient for very large or distributed settings.
\end{itemize}

\end{document}
